\documentclass[11pt]{article}

    \usepackage[T2A]{fontenc}
\usepackage[utf8]{inputenc}
\usepackage[english, russian]{babel}

\usepackage[breakable]{tcolorbox}
    \usepackage{parskip} % Stop auto-indenting (to mimic markdown behaviour)
    
    \usepackage{iftex}
    \ifPDFTeX
    	\usepackage[T1]{fontenc}
    	\usepackage{mathpazo}
    \else
    	\usepackage{fontspec}
    \fi

    % Basic figure setup, for now with no caption control since it's done
    % automatically by Pandoc (which extracts ![](path) syntax from Markdown).
    \usepackage{graphicx}
    % Maintain compatibility with old templates. Remove in nbconvert 6.0
    \let\Oldincludegraphics\includegraphics
    % Ensure that by default, figures have no caption (until we provide a
    % proper Figure object with a Caption API and a way to capture that
    % in the conversion process - todo).
    \usepackage{caption}
    \DeclareCaptionFormat{nocaption}{}
    \captionsetup{format=nocaption,aboveskip=0pt,belowskip=0pt}

    \usepackage{float}
    \floatplacement{figure}{H} % forces figures to be placed at the correct location
    \usepackage{xcolor} % Allow colors to be defined
    \usepackage{enumerate} % Needed for markdown enumerations to work
    \usepackage{geometry} % Used to adjust the document margins
    \usepackage{amsmath} % Equations
    \usepackage{amssymb} % Equations
    \usepackage{textcomp} % defines textquotesingle
    % Hack from http://tex.stackexchange.com/a/47451/13684:
    \AtBeginDocument{%
        \def\PYZsq{\textquotesingle}% Upright quotes in Pygmentized code
    }
    \usepackage{upquote} % Upright quotes for verbatim code
    \usepackage{eurosym} % defines \euro
    \usepackage[mathletters]{ucs} % Extended unicode (utf-8) support
    \usepackage{fancyvrb} % verbatim replacement that allows latex
    \usepackage{grffile} % extends the file name processing of package graphics 
                         % to support a larger range
    \makeatletter % fix for old versions of grffile with XeLaTeX
    \@ifpackagelater{grffile}{2019/11/01}
    {
      % Do nothing on new versions
    }
    {
      \def\Gread@@xetex#1{%
        \IfFileExists{"\Gin@base".bb}%
        {\Gread@eps{\Gin@base.bb}}%
        {\Gread@@xetex@aux#1}%
      }
    }
    \makeatother
    \usepackage[Export]{adjustbox} % Used to constrain images to a maximum size
    \adjustboxset{max size={0.9\linewidth}{0.9\paperheight}}

    % The hyperref package gives us a pdf with properly built
    % internal navigation ('pdf bookmarks' for the table of contents,
    % internal cross-reference links, web links for URLs, etc.)
    \usepackage{hyperref}
    % The default LaTeX title has an obnoxious amount of whitespace. By default,
    % titling removes some of it. It also provides customization options.
    \usepackage{titling}
    \usepackage{longtable} % longtable support required by pandoc >1.10
    \usepackage{booktabs}  % table support for pandoc > 1.12.2
    \usepackage[inline]{enumitem} % IRkernel/repr support (it uses the enumerate* environment)
    \usepackage[normalem]{ulem} % ulem is needed to support strikethroughs (\sout)
                                % normalem makes italics be italics, not underlines
    \usepackage{mathrsfs}
    
\usepackage{wrapfig}
\usepackage[rightcaption]{sidecap}
\providecommand{\keywords}[1]{\textbf{\textit{Keywords:}} #1}

\author{A. Yu. Drozdov}

    
    % Colors for the hyperref package
    \definecolor{urlcolor}{rgb}{0,.145,.698}
    \definecolor{linkcolor}{rgb}{.71,0.21,0.01}
    \definecolor{citecolor}{rgb}{.12,.54,.11}

    % ANSI colors
    \definecolor{ansi-black}{HTML}{3E424D}
    \definecolor{ansi-black-intense}{HTML}{282C36}
    \definecolor{ansi-red}{HTML}{E75C58}
    \definecolor{ansi-red-intense}{HTML}{B22B31}
    \definecolor{ansi-green}{HTML}{00A250}
    \definecolor{ansi-green-intense}{HTML}{007427}
    \definecolor{ansi-yellow}{HTML}{DDB62B}
    \definecolor{ansi-yellow-intense}{HTML}{B27D12}
    \definecolor{ansi-blue}{HTML}{208FFB}
    \definecolor{ansi-blue-intense}{HTML}{0065CA}
    \definecolor{ansi-magenta}{HTML}{D160C4}
    \definecolor{ansi-magenta-intense}{HTML}{A03196}
    \definecolor{ansi-cyan}{HTML}{60C6C8}
    \definecolor{ansi-cyan-intense}{HTML}{258F8F}
    \definecolor{ansi-white}{HTML}{C5C1B4}
    \definecolor{ansi-white-intense}{HTML}{A1A6B2}
    \definecolor{ansi-default-inverse-fg}{HTML}{FFFFFF}
    \definecolor{ansi-default-inverse-bg}{HTML}{000000}

    % common color for the border for error outputs.
    \definecolor{outerrorbackground}{HTML}{FFDFDF}

    % commands and environments needed by pandoc snippets
    % extracted from the output of `pandoc -s`
    \providecommand{\tightlist}{%
      \setlength{\itemsep}{0pt}\setlength{\parskip}{0pt}}
    \DefineVerbatimEnvironment{Highlighting}{Verbatim}{commandchars=\\\{\}}
    % Add ',fontsize=\small' for more characters per line
    \newenvironment{Shaded}{}{}
    \newcommand{\KeywordTok}[1]{\textcolor[rgb]{0.00,0.44,0.13}{\textbf{{#1}}}}
    \newcommand{\DataTypeTok}[1]{\textcolor[rgb]{0.56,0.13,0.00}{{#1}}}
    \newcommand{\DecValTok}[1]{\textcolor[rgb]{0.25,0.63,0.44}{{#1}}}
    \newcommand{\BaseNTok}[1]{\textcolor[rgb]{0.25,0.63,0.44}{{#1}}}
    \newcommand{\FloatTok}[1]{\textcolor[rgb]{0.25,0.63,0.44}{{#1}}}
    \newcommand{\CharTok}[1]{\textcolor[rgb]{0.25,0.44,0.63}{{#1}}}
    \newcommand{\StringTok}[1]{\textcolor[rgb]{0.25,0.44,0.63}{{#1}}}
    \newcommand{\CommentTok}[1]{\textcolor[rgb]{0.38,0.63,0.69}{\textit{{#1}}}}
    \newcommand{\OtherTok}[1]{\textcolor[rgb]{0.00,0.44,0.13}{{#1}}}
    \newcommand{\AlertTok}[1]{\textcolor[rgb]{1.00,0.00,0.00}{\textbf{{#1}}}}
    \newcommand{\FunctionTok}[1]{\textcolor[rgb]{0.02,0.16,0.49}{{#1}}}
    \newcommand{\RegionMarkerTok}[1]{{#1}}
    \newcommand{\ErrorTok}[1]{\textcolor[rgb]{1.00,0.00,0.00}{\textbf{{#1}}}}
    \newcommand{\NormalTok}[1]{{#1}}
    
    % Additional commands for more recent versions of Pandoc
    \newcommand{\ConstantTok}[1]{\textcolor[rgb]{0.53,0.00,0.00}{{#1}}}
    \newcommand{\SpecialCharTok}[1]{\textcolor[rgb]{0.25,0.44,0.63}{{#1}}}
    \newcommand{\VerbatimStringTok}[1]{\textcolor[rgb]{0.25,0.44,0.63}{{#1}}}
    \newcommand{\SpecialStringTok}[1]{\textcolor[rgb]{0.73,0.40,0.53}{{#1}}}
    \newcommand{\ImportTok}[1]{{#1}}
    \newcommand{\DocumentationTok}[1]{\textcolor[rgb]{0.73,0.13,0.13}{\textit{{#1}}}}
    \newcommand{\AnnotationTok}[1]{\textcolor[rgb]{0.38,0.63,0.69}{\textbf{\textit{{#1}}}}}
    \newcommand{\CommentVarTok}[1]{\textcolor[rgb]{0.38,0.63,0.69}{\textbf{\textit{{#1}}}}}
    \newcommand{\VariableTok}[1]{\textcolor[rgb]{0.10,0.09,0.49}{{#1}}}
    \newcommand{\ControlFlowTok}[1]{\textcolor[rgb]{0.00,0.44,0.13}{\textbf{{#1}}}}
    \newcommand{\OperatorTok}[1]{\textcolor[rgb]{0.40,0.40,0.40}{{#1}}}
    \newcommand{\BuiltInTok}[1]{{#1}}
    \newcommand{\ExtensionTok}[1]{{#1}}
    \newcommand{\PreprocessorTok}[1]{\textcolor[rgb]{0.74,0.48,0.00}{{#1}}}
    \newcommand{\AttributeTok}[1]{\textcolor[rgb]{0.49,0.56,0.16}{{#1}}}
    \newcommand{\InformationTok}[1]{\textcolor[rgb]{0.38,0.63,0.69}{\textbf{\textit{{#1}}}}}
    \newcommand{\WarningTok}[1]{\textcolor[rgb]{0.38,0.63,0.69}{\textbf{\textit{{#1}}}}}
    
    
    % Define a nice break command that doesn't care if a line doesn't already
    % exist.
    \def\br{\hspace*{\fill} \\* }
    % Math Jax compatibility definitions
    \def\gt{>}
    \def\lt{<}
    \let\Oldtex\TeX
    \let\Oldlatex\LaTeX
    \renewcommand{\TeX}{\textrm{\Oldtex}}
    \renewcommand{\LaTeX}{\textrm{\Oldlatex}}
    % Document parameters
    % Document title
    \title{04.do\_Lienard-Wichert\_potentials\_correspond\_to\_Lorentz\_calibration}
    
    
    
    
    
% Pygments definitions
\makeatletter
\def\PY@reset{\let\PY@it=\relax \let\PY@bf=\relax%
    \let\PY@ul=\relax \let\PY@tc=\relax%
    \let\PY@bc=\relax \let\PY@ff=\relax}
\def\PY@tok#1{\csname PY@tok@#1\endcsname}
\def\PY@toks#1+{\ifx\relax#1\empty\else%
    \PY@tok{#1}\expandafter\PY@toks\fi}
\def\PY@do#1{\PY@bc{\PY@tc{\PY@ul{%
    \PY@it{\PY@bf{\PY@ff{#1}}}}}}}
\def\PY#1#2{\PY@reset\PY@toks#1+\relax+\PY@do{#2}}

\@namedef{PY@tok@w}{\def\PY@tc##1{\textcolor[rgb]{0.73,0.73,0.73}{##1}}}
\@namedef{PY@tok@c}{\let\PY@it=\textit\def\PY@tc##1{\textcolor[rgb]{0.25,0.50,0.50}{##1}}}
\@namedef{PY@tok@cp}{\def\PY@tc##1{\textcolor[rgb]{0.74,0.48,0.00}{##1}}}
\@namedef{PY@tok@k}{\let\PY@bf=\textbf\def\PY@tc##1{\textcolor[rgb]{0.00,0.50,0.00}{##1}}}
\@namedef{PY@tok@kp}{\def\PY@tc##1{\textcolor[rgb]{0.00,0.50,0.00}{##1}}}
\@namedef{PY@tok@kt}{\def\PY@tc##1{\textcolor[rgb]{0.69,0.00,0.25}{##1}}}
\@namedef{PY@tok@o}{\def\PY@tc##1{\textcolor[rgb]{0.40,0.40,0.40}{##1}}}
\@namedef{PY@tok@ow}{\let\PY@bf=\textbf\def\PY@tc##1{\textcolor[rgb]{0.67,0.13,1.00}{##1}}}
\@namedef{PY@tok@nb}{\def\PY@tc##1{\textcolor[rgb]{0.00,0.50,0.00}{##1}}}
\@namedef{PY@tok@nf}{\def\PY@tc##1{\textcolor[rgb]{0.00,0.00,1.00}{##1}}}
\@namedef{PY@tok@nc}{\let\PY@bf=\textbf\def\PY@tc##1{\textcolor[rgb]{0.00,0.00,1.00}{##1}}}
\@namedef{PY@tok@nn}{\let\PY@bf=\textbf\def\PY@tc##1{\textcolor[rgb]{0.00,0.00,1.00}{##1}}}
\@namedef{PY@tok@ne}{\let\PY@bf=\textbf\def\PY@tc##1{\textcolor[rgb]{0.82,0.25,0.23}{##1}}}
\@namedef{PY@tok@nv}{\def\PY@tc##1{\textcolor[rgb]{0.10,0.09,0.49}{##1}}}
\@namedef{PY@tok@no}{\def\PY@tc##1{\textcolor[rgb]{0.53,0.00,0.00}{##1}}}
\@namedef{PY@tok@nl}{\def\PY@tc##1{\textcolor[rgb]{0.63,0.63,0.00}{##1}}}
\@namedef{PY@tok@ni}{\let\PY@bf=\textbf\def\PY@tc##1{\textcolor[rgb]{0.60,0.60,0.60}{##1}}}
\@namedef{PY@tok@na}{\def\PY@tc##1{\textcolor[rgb]{0.49,0.56,0.16}{##1}}}
\@namedef{PY@tok@nt}{\let\PY@bf=\textbf\def\PY@tc##1{\textcolor[rgb]{0.00,0.50,0.00}{##1}}}
\@namedef{PY@tok@nd}{\def\PY@tc##1{\textcolor[rgb]{0.67,0.13,1.00}{##1}}}
\@namedef{PY@tok@s}{\def\PY@tc##1{\textcolor[rgb]{0.73,0.13,0.13}{##1}}}
\@namedef{PY@tok@sd}{\let\PY@it=\textit\def\PY@tc##1{\textcolor[rgb]{0.73,0.13,0.13}{##1}}}
\@namedef{PY@tok@si}{\let\PY@bf=\textbf\def\PY@tc##1{\textcolor[rgb]{0.73,0.40,0.53}{##1}}}
\@namedef{PY@tok@se}{\let\PY@bf=\textbf\def\PY@tc##1{\textcolor[rgb]{0.73,0.40,0.13}{##1}}}
\@namedef{PY@tok@sr}{\def\PY@tc##1{\textcolor[rgb]{0.73,0.40,0.53}{##1}}}
\@namedef{PY@tok@ss}{\def\PY@tc##1{\textcolor[rgb]{0.10,0.09,0.49}{##1}}}
\@namedef{PY@tok@sx}{\def\PY@tc##1{\textcolor[rgb]{0.00,0.50,0.00}{##1}}}
\@namedef{PY@tok@m}{\def\PY@tc##1{\textcolor[rgb]{0.40,0.40,0.40}{##1}}}
\@namedef{PY@tok@gh}{\let\PY@bf=\textbf\def\PY@tc##1{\textcolor[rgb]{0.00,0.00,0.50}{##1}}}
\@namedef{PY@tok@gu}{\let\PY@bf=\textbf\def\PY@tc##1{\textcolor[rgb]{0.50,0.00,0.50}{##1}}}
\@namedef{PY@tok@gd}{\def\PY@tc##1{\textcolor[rgb]{0.63,0.00,0.00}{##1}}}
\@namedef{PY@tok@gi}{\def\PY@tc##1{\textcolor[rgb]{0.00,0.63,0.00}{##1}}}
\@namedef{PY@tok@gr}{\def\PY@tc##1{\textcolor[rgb]{1.00,0.00,0.00}{##1}}}
\@namedef{PY@tok@ge}{\let\PY@it=\textit}
\@namedef{PY@tok@gs}{\let\PY@bf=\textbf}
\@namedef{PY@tok@gp}{\let\PY@bf=\textbf\def\PY@tc##1{\textcolor[rgb]{0.00,0.00,0.50}{##1}}}
\@namedef{PY@tok@go}{\def\PY@tc##1{\textcolor[rgb]{0.53,0.53,0.53}{##1}}}
\@namedef{PY@tok@gt}{\def\PY@tc##1{\textcolor[rgb]{0.00,0.27,0.87}{##1}}}
\@namedef{PY@tok@err}{\def\PY@bc##1{{\setlength{\fboxsep}{\string -\fboxrule}\fcolorbox[rgb]{1.00,0.00,0.00}{1,1,1}{\strut ##1}}}}
\@namedef{PY@tok@kc}{\let\PY@bf=\textbf\def\PY@tc##1{\textcolor[rgb]{0.00,0.50,0.00}{##1}}}
\@namedef{PY@tok@kd}{\let\PY@bf=\textbf\def\PY@tc##1{\textcolor[rgb]{0.00,0.50,0.00}{##1}}}
\@namedef{PY@tok@kn}{\let\PY@bf=\textbf\def\PY@tc##1{\textcolor[rgb]{0.00,0.50,0.00}{##1}}}
\@namedef{PY@tok@kr}{\let\PY@bf=\textbf\def\PY@tc##1{\textcolor[rgb]{0.00,0.50,0.00}{##1}}}
\@namedef{PY@tok@bp}{\def\PY@tc##1{\textcolor[rgb]{0.00,0.50,0.00}{##1}}}
\@namedef{PY@tok@fm}{\def\PY@tc##1{\textcolor[rgb]{0.00,0.00,1.00}{##1}}}
\@namedef{PY@tok@vc}{\def\PY@tc##1{\textcolor[rgb]{0.10,0.09,0.49}{##1}}}
\@namedef{PY@tok@vg}{\def\PY@tc##1{\textcolor[rgb]{0.10,0.09,0.49}{##1}}}
\@namedef{PY@tok@vi}{\def\PY@tc##1{\textcolor[rgb]{0.10,0.09,0.49}{##1}}}
\@namedef{PY@tok@vm}{\def\PY@tc##1{\textcolor[rgb]{0.10,0.09,0.49}{##1}}}
\@namedef{PY@tok@sa}{\def\PY@tc##1{\textcolor[rgb]{0.73,0.13,0.13}{##1}}}
\@namedef{PY@tok@sb}{\def\PY@tc##1{\textcolor[rgb]{0.73,0.13,0.13}{##1}}}
\@namedef{PY@tok@sc}{\def\PY@tc##1{\textcolor[rgb]{0.73,0.13,0.13}{##1}}}
\@namedef{PY@tok@dl}{\def\PY@tc##1{\textcolor[rgb]{0.73,0.13,0.13}{##1}}}
\@namedef{PY@tok@s2}{\def\PY@tc##1{\textcolor[rgb]{0.73,0.13,0.13}{##1}}}
\@namedef{PY@tok@sh}{\def\PY@tc##1{\textcolor[rgb]{0.73,0.13,0.13}{##1}}}
\@namedef{PY@tok@s1}{\def\PY@tc##1{\textcolor[rgb]{0.73,0.13,0.13}{##1}}}
\@namedef{PY@tok@mb}{\def\PY@tc##1{\textcolor[rgb]{0.40,0.40,0.40}{##1}}}
\@namedef{PY@tok@mf}{\def\PY@tc##1{\textcolor[rgb]{0.40,0.40,0.40}{##1}}}
\@namedef{PY@tok@mh}{\def\PY@tc##1{\textcolor[rgb]{0.40,0.40,0.40}{##1}}}
\@namedef{PY@tok@mi}{\def\PY@tc##1{\textcolor[rgb]{0.40,0.40,0.40}{##1}}}
\@namedef{PY@tok@il}{\def\PY@tc##1{\textcolor[rgb]{0.40,0.40,0.40}{##1}}}
\@namedef{PY@tok@mo}{\def\PY@tc##1{\textcolor[rgb]{0.40,0.40,0.40}{##1}}}
\@namedef{PY@tok@ch}{\let\PY@it=\textit\def\PY@tc##1{\textcolor[rgb]{0.25,0.50,0.50}{##1}}}
\@namedef{PY@tok@cm}{\let\PY@it=\textit\def\PY@tc##1{\textcolor[rgb]{0.25,0.50,0.50}{##1}}}
\@namedef{PY@tok@cpf}{\let\PY@it=\textit\def\PY@tc##1{\textcolor[rgb]{0.25,0.50,0.50}{##1}}}
\@namedef{PY@tok@c1}{\let\PY@it=\textit\def\PY@tc##1{\textcolor[rgb]{0.25,0.50,0.50}{##1}}}
\@namedef{PY@tok@cs}{\let\PY@it=\textit\def\PY@tc##1{\textcolor[rgb]{0.25,0.50,0.50}{##1}}}

\def\PYZbs{\char`\\}
\def\PYZus{\char`\_}
\def\PYZob{\char`\{}
\def\PYZcb{\char`\}}
\def\PYZca{\char`\^}
\def\PYZam{\char`\&}
\def\PYZlt{\char`\<}
\def\PYZgt{\char`\>}
\def\PYZsh{\char`\#}
\def\PYZpc{\char`\%}
\def\PYZdl{\char`\$}
\def\PYZhy{\char`\-}
\def\PYZsq{\char`\'}
\def\PYZdq{\char`\"}
\def\PYZti{\char`\~}
% for compatibility with earlier versions
\def\PYZat{@}
\def\PYZlb{[}
\def\PYZrb{]}
\makeatother


    % For linebreaks inside Verbatim environment from package fancyvrb. 
    \makeatletter
        \newbox\Wrappedcontinuationbox 
        \newbox\Wrappedvisiblespacebox 
        \newcommand*\Wrappedvisiblespace {\textcolor{red}{\textvisiblespace}} 
        \newcommand*\Wrappedcontinuationsymbol {\textcolor{red}{\llap{\tiny$\m@th\hookrightarrow$}}} 
        \newcommand*\Wrappedcontinuationindent {3ex } 
        \newcommand*\Wrappedafterbreak {\kern\Wrappedcontinuationindent\copy\Wrappedcontinuationbox} 
        % Take advantage of the already applied Pygments mark-up to insert 
        % potential linebreaks for TeX processing. 
        %        {, <, #, %, $, ' and ": go to next line. 
        %        _, }, ^, &, >, - and ~: stay at end of broken line. 
        % Use of \textquotesingle for straight quote. 
        \newcommand*\Wrappedbreaksatspecials {% 
            \def\PYGZus{\discretionary{\char`\_}{\Wrappedafterbreak}{\char`\_}}% 
            \def\PYGZob{\discretionary{}{\Wrappedafterbreak\char`\{}{\char`\{}}% 
            \def\PYGZcb{\discretionary{\char`\}}{\Wrappedafterbreak}{\char`\}}}% 
            \def\PYGZca{\discretionary{\char`\^}{\Wrappedafterbreak}{\char`\^}}% 
            \def\PYGZam{\discretionary{\char`\&}{\Wrappedafterbreak}{\char`\&}}% 
            \def\PYGZlt{\discretionary{}{\Wrappedafterbreak\char`\<}{\char`\<}}% 
            \def\PYGZgt{\discretionary{\char`\>}{\Wrappedafterbreak}{\char`\>}}% 
            \def\PYGZsh{\discretionary{}{\Wrappedafterbreak\char`\#}{\char`\#}}% 
            \def\PYGZpc{\discretionary{}{\Wrappedafterbreak\char`\%}{\char`\%}}% 
            \def\PYGZdl{\discretionary{}{\Wrappedafterbreak\char`\$}{\char`\$}}% 
            \def\PYGZhy{\discretionary{\char`\-}{\Wrappedafterbreak}{\char`\-}}% 
            \def\PYGZsq{\discretionary{}{\Wrappedafterbreak\textquotesingle}{\textquotesingle}}% 
            \def\PYGZdq{\discretionary{}{\Wrappedafterbreak\char`\"}{\char`\"}}% 
            \def\PYGZti{\discretionary{\char`\~}{\Wrappedafterbreak}{\char`\~}}% 
        } 
        % Some characters . , ; ? ! / are not pygmentized. 
        % This macro makes them "active" and they will insert potential linebreaks 
        \newcommand*\Wrappedbreaksatpunct {% 
            \lccode`\~`\.\lowercase{\def~}{\discretionary{\hbox{\char`\.}}{\Wrappedafterbreak}{\hbox{\char`\.}}}% 
            \lccode`\~`\,\lowercase{\def~}{\discretionary{\hbox{\char`\,}}{\Wrappedafterbreak}{\hbox{\char`\,}}}% 
            \lccode`\~`\;\lowercase{\def~}{\discretionary{\hbox{\char`\;}}{\Wrappedafterbreak}{\hbox{\char`\;}}}% 
            \lccode`\~`\:\lowercase{\def~}{\discretionary{\hbox{\char`\:}}{\Wrappedafterbreak}{\hbox{\char`\:}}}% 
            \lccode`\~`\?\lowercase{\def~}{\discretionary{\hbox{\char`\?}}{\Wrappedafterbreak}{\hbox{\char`\?}}}% 
            \lccode`\~`\!\lowercase{\def~}{\discretionary{\hbox{\char`\!}}{\Wrappedafterbreak}{\hbox{\char`\!}}}% 
            \lccode`\~`\/\lowercase{\def~}{\discretionary{\hbox{\char`\/}}{\Wrappedafterbreak}{\hbox{\char`\/}}}% 
            \catcode`\.\active
            \catcode`\,\active 
            \catcode`\;\active
            \catcode`\:\active
            \catcode`\?\active
            \catcode`\!\active
            \catcode`\/\active 
            \lccode`\~`\~ 	
        }
    \makeatother

    \let\OriginalVerbatim=\Verbatim
    \makeatletter
    \renewcommand{\Verbatim}[1][1]{%
        %\parskip\z@skip
        \sbox\Wrappedcontinuationbox {\Wrappedcontinuationsymbol}%
        \sbox\Wrappedvisiblespacebox {\FV@SetupFont\Wrappedvisiblespace}%
        \def\FancyVerbFormatLine ##1{\hsize\linewidth
            \vtop{\raggedright\hyphenpenalty\z@\exhyphenpenalty\z@
                \doublehyphendemerits\z@\finalhyphendemerits\z@
                \strut ##1\strut}%
        }%
        % If the linebreak is at a space, the latter will be displayed as visible
        % space at end of first line, and a continuation symbol starts next line.
        % Stretch/shrink are however usually zero for typewriter font.
        \def\FV@Space {%
            \nobreak\hskip\z@ plus\fontdimen3\font minus\fontdimen4\font
            \discretionary{\copy\Wrappedvisiblespacebox}{\Wrappedafterbreak}
            {\kern\fontdimen2\font}%
        }%
        
        % Allow breaks at special characters using \PYG... macros.
        \Wrappedbreaksatspecials
        % Breaks at punctuation characters . , ; ? ! and / need catcode=\active 	
        \OriginalVerbatim[#1,codes*=\Wrappedbreaksatpunct]%
    }
    \makeatother

    % Exact colors from NB
    \definecolor{incolor}{HTML}{303F9F}
    \definecolor{outcolor}{HTML}{D84315}
    \definecolor{cellborder}{HTML}{CFCFCF}
    \definecolor{cellbackground}{HTML}{F7F7F7}
    
    % prompt
    \makeatletter
    \newcommand{\boxspacing}{\kern\kvtcb@left@rule\kern\kvtcb@boxsep}
    \makeatother
    \newcommand{\prompt}[4]{
        {\ttfamily\llap{{\color{#2}[#3]:\hspace{3pt}#4}}\vspace{-\baselineskip}}
    }
    

    
    % Prevent overflowing lines due to hard-to-break entities
    \sloppy 
    % Setup hyperref package
    \hypersetup{
      breaklinks=true,  % so long urls are correctly broken across lines
      colorlinks=true,
      urlcolor=urlcolor,
      linkcolor=linkcolor,
      citecolor=citecolor,
      }
    % Slightly bigger margins than the latex defaults
    
    \geometry{verbose,tmargin=1in,bmargin=1in,lmargin=1in,rmargin=1in}
    
    

\begin{document}
    
    \maketitle
    
    

    
    \section{Соответствуют ли потенциалы Лиенара-Вихерта калибровке
Лоренца?}\label{ux441ux43eux43eux442ux432ux435ux442ux441ux442ux432ux443ux44eux442-ux43bux438-ux43fux43eux442ux435ux43dux446ux438ux430ux43bux44b-ux43bux438ux435ux43dux430ux440ux430-ux432ux438ux445ux435ux440ux442ux430-ux43aux430ux43bux438ux431ux440ux43eux432ux43aux435-ux43bux43eux440ux435ux43dux446ux430}

Чтобы проверить сомнение о применимости калибровки Лоренца к
классическому векторному потенциалу в электродинамике, удобно
воспользоваться таким математическим инструментом, как потенциалы
Лиенара-Вихерта, которые определены как потенциалы поля, создаваемого
движущимся одиночным (точечным) зарядом. Потенциалы Лиенара-Вихерта, как
известно {[}CITATION ЕМЛ73 \l 1049 {]} получаются в результате
интегрирования запаздывающих потенциалов 12 и 13, которые в свою очередь
являются решением волновых уравнений Даламбера, которые получены из
уравнений Максвелла с применением калибровки Лоренца, хотя и не вполне
корректным путём в том смысле, что там, где авторам было удобно
применялась калибровка Лоренца, а там, где не удобно, - неявно
применялась калибровка Кулона из электростатики. Итак, записываем
потенциалы ЛВ
\[\varphi =\frac{1}{\varepsilon }\frac{e}{\left( R-\frac{\overrightarrow{v}\left( x',y',z',t' \right)\cdot \overrightarrow{R}}{c} \right)}\]
\[\overrightarrow{{{A}_{H}}}=\frac{1}{c}\frac{e\ \overrightarrow{v}\left( x',y',z',t' \right)}{\ R-\frac{\overrightarrow{v}\left( x',y',z',t' \right)\cdot \overrightarrow{R}}{c}}\]
Производная скалярного потенциала по времени равна
\[\frac{\partial \varphi }{\partial t}=\frac{\partial \varphi }{\partial t'}\frac{\partial t'}{\partial t}\]
Найдём производную запаздывающего момента \(t'=t-\frac{R}{c}\) по
времени \[\frac{\partial t'}{\partial t}\] Дифференцируя
\[R\left( \overrightarrow{x},\ \overrightarrow{x'}\left( t' \right) \right)=c\left( t-t' \right)\]
по \(t\) получаем
\(\frac{\partial R}{\partial t}=\frac{\partial R}{\partial t'}\frac{\partial t'}{\partial t}=c\left( 1-\frac{\partial t'}{\partial t} \right)\).
Значение \(\frac{\partial R}{\partial t'}\)находим так
\[\frac{\partial R\left( t' \right)}{\partial t'}=\frac{\partial }{\partial t'}{{\left( \overrightarrow{R}\cdot \overrightarrow{R} \right)}^{{}^{1}/{}_{2}}}=\frac{1}{2}{{\left( \overrightarrow{R}\cdot \overrightarrow{R} \right)}^{-{}^{1}/{}_{2}}}\frac{\partial }{\partial t'}\left( \overrightarrow{R}\cdot \overrightarrow{R} \right)=\frac{1}{2R}\left( \overrightarrow{R}\cdot \frac{\partial \overrightarrow{R}}{\partial t'}+\frac{\partial \overrightarrow{R}}{\partial t'}\cdot \overrightarrow{R} \right)=\frac{\overrightarrow{R}\left( t' \right)}{R\left( t' \right)}\cdot \frac{\partial \overrightarrow{R}\left( t' \right)}{\partial t'}=-\frac{\overrightarrow{R}\cdot \overrightarrow{v}}{R}\]
Далее учитывая что
\(\frac{\partial \overrightarrow{R}}{\partial t'}=-\overrightarrow{v}\left( t' \right)\)
имеем
\[\frac{\partial R}{\partial t'}=-\frac{\overrightarrow{R}}{R}\cdot \overrightarrow{v}\]
теперь
\(\frac{\partial R}{\partial t'}\frac{\partial t'}{\partial t}=c\left( 1-\frac{\partial t'}{\partial t} \right)\)преобразуется
к
\(-\frac{\overrightarrow{R}\cdot \overrightarrow{v}}{cR}\frac{\partial t'}{\partial t}=1-\frac{\partial t'}{\partial t}\)
решая которое получаем
\[\frac{\partial t'}{\partial t}=\frac{1}{1-\frac{\overrightarrow{v}\cdot \overrightarrow{R}}{Rc}}\]
Для производной скалярного потенциала $\varphi$ по \(t'\) вводим
обозначение \(K=R-\frac{\overrightarrow{v}\cdot \overrightarrow{R}}{c}\)
благодаря которому запись скалярного потенциала упрощается
\[\varphi =\frac{1}{\varepsilon }\frac{e}{K}\] кроме того
\[\frac{\partial t'}{\partial t}=\frac{R}{K}\]
\(\frac{\partial \varphi }{\partial t'}=\frac{\partial \varphi }{\partial K}\frac{\partial K}{\partial t'}=-\frac{e}{\varepsilon {{K}^{2}}}\left\{ \frac{\partial R}{\partial t'}-\frac{1}{c}\frac{\partial }{\partial t'}\left( \overrightarrow{v}\cdot \overrightarrow{R} \right) \right\}\)
\(\frac{\partial }{\partial t'}\left( \overrightarrow{v}\cdot \overrightarrow{R} \right)=\overrightarrow{v}\frac{\partial \overrightarrow{R}}{\partial t'}+\overrightarrow{R}\frac{\partial \overrightarrow{v}}{\partial t'}=-\overrightarrow{v}\cdot \overrightarrow{v}+\overrightarrow{R}\cdot \overrightarrow{a}\)
\(\frac{\partial \varphi }{\partial t'}=-\frac{e}{\varepsilon {{K}^{2}}}\left\{ -\frac{\overrightarrow{R}}{R}\cdot \overrightarrow{v}+\frac{1}{c}\left( \overrightarrow{v}\cdot \overrightarrow{v}-\overrightarrow{R}\cdot \overrightarrow{a} \right) \right\}\)
Итак
\[\frac{\partial \varphi }{\partial t}=-\frac{e}{\varepsilon }\frac{R\frac{\vec{v}\cdot \vec{v}-\overrightarrow{a}\cdot \overrightarrow{R}}{c}-~\overrightarrow{v}\cdot \overrightarrow{R}}{{{K}^{3}}}\]
Дивергенция векторного потенциала
\[\nabla \overrightarrow{{{A}_{H}}}=\frac{e}{c}\nabla \frac{\overrightarrow{v}}{K}=\frac{e}{c}\left\{ \frac{1}{K}\nabla \overrightarrow{v}+\overrightarrow{v}\cdot \nabla \frac{1}{K} \right\}\]
Дифференцируя скорость точечного заряда по координатам точки наблюдения
\[{{\nabla }_{a}}\ \overrightarrow{v}\left( x',y',z',t-\frac{R}{v} \right)=\frac{\partial {{v}_{x}}}{\partial t'}\frac{\partial t'}{\partial x}+\frac{\partial {{v}_{y}}}{\partial t'}\frac{\partial t'}{\partial y}+\frac{\partial {{v}_{z}}}{\partial t'}\frac{\partial t'}{\partial z}=\frac{\partial \overrightarrow{v}}{\partial t'}\cdot \nabla t'\]
Градиент запаздывающего момента времени \[\nabla t'\]
\[R\left( \overrightarrow{x},\ \overrightarrow{x'}\left( t' \right) \right)=c\left( t-t' \right)\]
\[t'=t-\frac{R\left( \overrightarrow{x},\overrightarrow{x'}\left( t' \right) \right)}{c}\]
Дифференцируем по координатам
\[\nabla t'=-\frac{1}{c}\nabla R\left( \overrightarrow{x},\overrightarrow{x'}\left( t' \right) \right)\]
\[\nabla R\left( \overrightarrow{x},\overrightarrow{x'}\left( t' \right) \right)=\nabla R{{\left( \overrightarrow{x},\overrightarrow{x'}\left( t' \right) \right)}_{t'=const}}+\nabla R{{\left( \overrightarrow{x},\overrightarrow{x'}\left( t' \right) \right)}_{\overrightarrow{x}=const}}=\frac{\overrightarrow{R}}{R}+\frac{\partial R}{\partial t'}\nabla t'\]
\[\nabla t'=-\frac{1}{c}\left( \frac{\overrightarrow{R}}{R}+\frac{\partial R}{\partial t'}\nabla t' \right)\]
Далее
\[\nabla t'=-\frac{1}{c}\left( \frac{\overrightarrow{R}}{R}+\frac{\overrightarrow{R}\left( t' \right)}{R\left( t' \right)}\cdot \frac{\partial \overrightarrow{R}\left( t' \right)}{\partial t'}\nabla t' \right)\]
Учитывая что
\[\frac{\partial \overrightarrow{R}\left( t' \right)}{\partial t'}=-\overrightarrow{v}\left( t' \right)\]
\[\nabla t'=-\frac{1}{c}\left( \frac{\overrightarrow{R}}{R}-\frac{\overrightarrow{R}\left( t' \right)}{R\left( t' \right)}\cdot \overrightarrow{v}\left( t' \right)\nabla t' \right)\]
\[\nabla t'=-\frac{1}{c}\frac{\overrightarrow{R}}{R}+\frac{1}{c}\frac{\overrightarrow{R}\left( t' \right)}{R\left( t' \right)}\cdot \overrightarrow{v}\left( t' \right)\nabla t'\]
\[\nabla t'\left( 1-\frac{1}{c}\frac{\overrightarrow{R}\left( t' \right)\cdot \overrightarrow{v}\left( t' \right)}{R\left( t' \right)} \right)=-\frac{1}{c}\frac{\overrightarrow{R}}{R}\]
\[\nabla t'=-\frac{\ \overrightarrow{R}}{c\ \left( R-\frac{\overrightarrow{v}\cdot \overrightarrow{R}}{c} \right)}=-\frac{\ \overrightarrow{R}}{c\ K}\]
Таким образом
\[\nabla \overrightarrow{v}=\frac{\partial \overrightarrow{v}}{\partial t'}\cdot \nabla t'=-\frac{\ \overrightarrow{a}\cdot \overrightarrow{R}}{c\ K}\]
Далее

\[\nabla \frac{1}{K}=-\frac{1}{{{K}^{2}}}\nabla K\]

\[\nabla K=\nabla \left( R-\frac{\overrightarrow{v}\cdot \overrightarrow{R}}{c} \right)=\nabla R-\frac{1}{c}\nabla \left( \overrightarrow{v}\cdot \overrightarrow{R} \right)\]

\[\begin{align}
& \nabla R=\frac{\overrightarrow{R}}{R}+\frac{\partial R}{\partial t'}\nabla t'=\frac{\overrightarrow{R}}{R}-\frac{\partial R}{\partial t'}\frac{\ \overrightarrow{R}}{c\ K}=\frac{\overrightarrow{R}}{R}-\left( \frac{\overrightarrow{R}}{R}\cdot \frac{\partial \overrightarrow{R}}{\partial t'} \right)\frac{\ \overrightarrow{R}}{c\ K}=\frac{\overrightarrow{R}}{R}+\left( \frac{\overrightarrow{R}\cdot \overrightarrow{v}}{R} \right)\frac{\ \overrightarrow{R}}{c\ K}=\frac{\overrightarrow{R}}{R}\left( 1+\frac{\overrightarrow{R}\cdot \overrightarrow{v}}{c\ K} \right) \\ 
& =\frac{\overrightarrow{R}}{R}\left( \frac{c\ K+\overrightarrow{R}\cdot \overrightarrow{v}}{c\ K} \right)=\frac{\overrightarrow{R}}{R}\left( \frac{c\ \left( R-\frac{\overrightarrow{R}\cdot \overrightarrow{v}}{c} \right)+\overrightarrow{R}\cdot \overrightarrow{v}}{c\ K} \right)=\frac{\overrightarrow{R}}{R}\left( \frac{c\ R-\overrightarrow{R}\cdot \overrightarrow{v}+\overrightarrow{R}\cdot \overrightarrow{v}}{c\ K} \right)=\frac{\overrightarrow{R}}{R}\frac{c\,R}{c\,K}=\frac{\overrightarrow{R}\ }{K}
\end{align}\]

\[\nabla \left( \overrightarrow{v}\cdot \overrightarrow{R} \right)=\overrightarrow{v}\times rot\overrightarrow{R}+\overrightarrow{R}\times rot\overrightarrow{v}+\left( \overrightarrow{v},\nabla  \right)\overrightarrow{R}+\left( \overrightarrow{R},\nabla  \right)\overrightarrow{v}\]

Учитывая, что в декартовой системе координат
\({{R}_{x}}=x-x'\left( t' \right)\) \({{R}_{y}}=y-y'\left( t' \right)\)
\({{R}_{x}}=z-z'\left( t' \right)\) Ищем выражение для ротора радиус
вектора \[rot\overrightarrow{R}=\left| \begin{matrix}
   \overrightarrow{i} & \overrightarrow{j} & \overrightarrow{k}  \\
   \frac{\partial }{\partial x} & \frac{\partial }{\partial y} & \frac{\partial }{\partial z}  \\
   {{R}_{x}} & {{R}_{y}} & {{R}_{z}}  \\
\end{matrix} \right|=\begin{matrix}
   \overrightarrow{i}\left( \frac{\partial {{R}_{z}}}{\partial y}-\frac{\partial {{R}_{y}}}{\partial z} \right)  \\
   \overrightarrow{j}\left( \frac{\partial {{R}_{x}}}{\partial z}-\frac{\partial {{R}_{z}}}{\partial x} \right)  \\
   \overrightarrow{k}\left( \frac{\partial {{R}_{y}}}{\partial x}-\frac{\partial {{R}_{x}}}{\partial y} \right)  \\
\end{matrix}=\begin{matrix}
   \overrightarrow{i}\left( -\frac{\partial z'}{\partial t'}\frac{\partial t'}{\partial y}+\frac{\partial y'}{\partial t'}\frac{\partial t'}{\partial z} \right)  \\
   \overrightarrow{j}\left( -\frac{\partial x'}{\partial t'}\frac{\partial t'}{\partial z}+\frac{\partial z'}{\partial t'}\frac{\partial t'}{\partial x} \right)  \\
   \overrightarrow{k}\left( -\frac{\partial y'}{\partial t'}\frac{\partial t'}{\partial x}+\frac{\partial x'}{\partial t'}\frac{\partial t'}{\partial y} \right)  \\
\end{matrix}=\begin{matrix}
   \frac{\overrightarrow{i}}{c}\left( {{v}_{z}}\frac{\partial R}{\partial y}-{{v}_{y}}\frac{\partial R}{\partial z} \right)  \\
   \frac{\overrightarrow{j}}{c}\left( {{v}_{x}}\frac{\partial R}{\partial z}-{{v}_{z}}\frac{\partial R}{\partial x} \right)  \\
   \frac{\overrightarrow{k}}{c}\left( {{v}_{y}}\frac{\partial R}{\partial x}-{{v}_{x}}\frac{\partial R}{\partial y} \right)  \\
\end{matrix}\] Производная \(R\) по координате точки наблюдения
\[\frac{\partial R}{\partial x}=\frac{\partial }{\partial x}{{\left( \overrightarrow{R}\cdot \overrightarrow{R} \right)}^{{}^{1}/{}_{2}}}=\frac{1}{2}{{\left( \overrightarrow{R}\cdot \overrightarrow{R} \right)}^{-{}^{1}/{}_{2}}}\frac{\partial }{\partial x}\left( \overrightarrow{R}\cdot \overrightarrow{R} \right)=\frac{1}{2R}\left( \overrightarrow{R}\cdot \frac{\partial \overrightarrow{R}}{\partial x}+\frac{\partial \overrightarrow{R}}{\partial x}\cdot \overrightarrow{R} \right)=\frac{\overrightarrow{R}}{R}\cdot \frac{\partial \overrightarrow{R}}{\partial x}\]
Поскольку в декартовой системе координат
\[\overrightarrow{R}=\overrightarrow{i}\ {{R}_{x}}\left( t' \right)+\overrightarrow{j}\ {{R}_{y}}\left( t' \right)+\overrightarrow{k}\ {{R}_{z}}\left( t' \right)=\overrightarrow{i}\left( x-x'\left( t' \right) \right)+\overrightarrow{j}\left( y-y'\left( t' \right) \right)+\overrightarrow{k}\left( z-z'\left( t' \right) \right)\]
постольку
\[\frac{\partial \overrightarrow{R}}{\partial x}=\begin{matrix}
   \overrightarrow{i}\ \frac{\partial {{R}_{x}}\left( t' \right)}{\partial x}  \\
   \overrightarrow{j}\ \frac{\partial {{R}_{y}}\left( t' \right)}{\partial x}  \\
   \overrightarrow{k}\ \frac{\partial {{R}_{z}}\left( t' \right)}{\partial x}  \\
\end{matrix}=\begin{matrix}
   \overrightarrow{i}\ \frac{\partial \left( x-x'\left( t' \right) \right)}{\partial x}  \\
   \overrightarrow{j}\ \frac{\partial \left( y-y'\left( t' \right) \right)}{\partial x}  \\
   \overrightarrow{k}\ \frac{\partial \left( z-z'\left( t' \right) \right)}{\partial x}  \\
\end{matrix}=\begin{matrix}
   \overrightarrow{i}\ \left( 1-\frac{\partial x'}{\partial t'}\frac{\partial t'}{\partial x} \right)  \\
   -\overrightarrow{j}\ \frac{\partial y'}{\partial t'}\frac{\partial t'}{\partial x}  \\
   -\overrightarrow{k}\ \frac{\partial z'}{\partial t'}\frac{\partial t'}{\partial x}  \\
\end{matrix}=\begin{matrix}
   \overrightarrow{i}\ \left( 1+\frac{{{v}_{x}}}{c}\frac{\partial R}{\partial x} \right)  \\
   \overrightarrow{j}\ \frac{{{v}_{y}}}{c}\frac{\partial R}{\partial x}  \\
   \overrightarrow{k}\ \frac{{{v}_{z}}}{c}\frac{\partial R}{\partial x}  \\
\end{matrix}=\overrightarrow{i}+\frac{\overrightarrow{v}}{c}\frac{\partial R}{\partial x}\]
Следовательно
\[\frac{\partial R}{\partial x}=\frac{\overrightarrow{R}}{R}\cdot \frac{\partial \overrightarrow{R}}{\partial x}=\frac{\overrightarrow{R}}{R}\cdot \left( \overrightarrow{i}+\frac{\overrightarrow{v}}{c}\frac{\partial R}{\partial x} \right)=\frac{{{R}_{x}}}{R}+\frac{\overrightarrow{R}\cdot \overrightarrow{v}}{Rc}\frac{\partial R}{\partial x}\]
\[\frac{\partial R}{\partial x}\left( 1-\frac{\overrightarrow{R}\cdot \overrightarrow{v}}{Rc} \right)=\frac{{{R}_{x}}}{R}\]
откуда
\[\frac{\partial R}{\partial x}=\frac{{{R}_{x}}}{R\left( 1-\frac{\overrightarrow{R}\cdot \overrightarrow{v}}{Rc} \right)}=\frac{{{R}_{x}}}{R-\frac{\overrightarrow{R}\cdot \overrightarrow{v}}{c}}=\frac{{{R}_{x}}}{K}\]
И
\[\frac{\partial \overrightarrow{R}}{\partial x}=\overrightarrow{i}+\frac{\overrightarrow{v}}{c}\frac{\partial R}{\partial x}=\overrightarrow{i}+\frac{\overrightarrow{v}}{c}\frac{{{R}_{x}}}{K}\]
Подставляя в выражение для ротора радиус вектора
\[rot\overrightarrow{R}=\begin{matrix}
   \frac{\overrightarrow{i}}{c}\left( {{v}_{z}}\frac{\partial R}{\partial y}-{{v}_{y}}\frac{\partial R}{\partial z} \right)  \\
   \frac{\overrightarrow{j}}{c}\left( {{v}_{x}}\frac{\partial R}{\partial z}-{{v}_{z}}\frac{\partial R}{\partial x} \right)  \\
   \frac{\overrightarrow{k}}{c}\left( {{v}_{y}}\frac{\partial R}{\partial x}-{{v}_{x}}\frac{\partial R}{\partial y} \right)  \\
\end{matrix}=\begin{matrix}
   \frac{\overrightarrow{i}}{cK}\left( {{v}_{z}}{{R}_{y}}-{{v}_{y}}{{R}_{z}} \right)  \\
   \frac{\overrightarrow{j}}{cK}\left( {{v}_{x}}{{R}_{z}}-{{v}_{z}}{{R}_{x}} \right)  \\
   \frac{\overrightarrow{k}}{cK}\left( {{v}_{y}}{{R}_{x}}-{{v}_{x}}{{R}_{y}} \right)  \\
\end{matrix}\] Теперь находим векторное произведение запаздывающей
скорости на ротор радиус вектора
\[cK\left( \overrightarrow{v}\times rot\overrightarrow{R} \right)=\left| \begin{matrix}
   \overrightarrow{i} & \overrightarrow{j} & \overrightarrow{k}  \\
   {{v}_{x}} & {{v}_{y}} & {{v}_{z}}  \\
   {{v}_{z}}{{R}_{y}}-{{v}_{y}}{{R}_{z}} & {{v}_{x}}{{R}_{z}}-{{v}_{z}}{{R}_{x}} & {{v}_{y}}{{R}_{x}}-{{v}_{x}}{{R}_{y}}  \\
\end{matrix} \right|=\begin{matrix}
   \overrightarrow{i}\left( {{v}_{y}}\left( {{v}_{y}}{{R}_{x}}-{{v}_{x}}{{R}_{y}} \right)-{{v}_{z}}\left( {{v}_{x}}{{R}_{z}}-{{v}_{z}}{{R}_{x}} \right) \right)  \\
   \overrightarrow{j}\left( {{v}_{z}}\left( {{v}_{z}}{{R}_{y}}-{{v}_{y}}{{R}_{z}} \right)-{{v}_{x}}\left( {{v}_{y}}{{R}_{x}}-{{v}_{x}}{{R}_{y}} \right) \right)  \\
   \overrightarrow{k}\left( {{v}_{x}}\left( {{v}_{x}}{{R}_{z}}-{{v}_{z}}{{R}_{x}} \right)-{{v}_{y}}\left( {{v}_{z}}{{R}_{y}}-{{v}_{y}}{{R}_{z}} \right) \right)  \\
\end{matrix}\] \[=\begin{matrix}
   \overrightarrow{i}\left( {{R}_{x}}\left( {{v}_{y}}^{2}+{{v}_{z}}^{2} \right)-{{v}_{x}}{{v}_{y}}{{R}_{y}}-{{v}_{x}}{{v}_{z}}{{R}_{z}} \right)  \\
   \overrightarrow{j}\left( {{R}_{y}}\left( {{v}_{x}}^{2}+{{v}_{z}}^{2} \right)-{{v}_{y}}{{v}_{x}}{{R}_{x}}-{{v}_{y}}{{v}_{z}}{{R}_{z}} \right)  \\
   \overrightarrow{k}\left( {{R}_{z}}\left( {{v}_{x}}^{2}+{{v}_{y}}^{2} \right)-{{v}_{z}}{{v}_{x}}{{R}_{x}}-{{v}_{z}}{{v}_{y}}{{R}_{y}} \right)  \\
\end{matrix}=\begin{matrix}
   \overrightarrow{i}\left( {{R}_{x}}\left( {{v}^{2}}-{{v}_{x}}^{2} \right)-{{v}_{x}}\left( {{v}_{y}}{{R}_{y}}-{{v}_{z}}{{R}_{z}} \right) \right)  \\
   \overrightarrow{j}\left( {{R}_{y}}\left( {{v}^{2}}-{{v}_{y}}^{2} \right)-{{v}_{y}}\left( {{v}_{x}}{{R}_{x}}-{{v}_{z}}{{R}_{z}} \right) \right)  \\
   \overrightarrow{k}\left( {{R}_{z}}\left( {{v}^{2}}-{{v}_{z}}^{2} \right)-{{v}_{z}}\left( {{v}_{x}}{{R}_{x}}-{{v}_{y}}{{R}_{y}} \right) \right)  \\
\end{matrix}=\begin{matrix}
   \overrightarrow{i}\left( {{R}_{x}}{{v}^{2}}-{{v}_{x}}\left( \overrightarrow{v}\cdot \overrightarrow{R} \right) \right)  \\
   \overrightarrow{j}\left( {{R}_{y}}{{v}^{2}}-{{v}_{y}}\left( \overrightarrow{v}\cdot \overrightarrow{R} \right) \right)  \\
   \overrightarrow{k}\left( {{R}_{z}}{{v}^{2}}-{{v}_{z}}\left( \overrightarrow{v}\cdot \overrightarrow{R} \right) \right)  \\
\end{matrix}\]

\[\overrightarrow{v}\times rot\overrightarrow{R}=\frac{{{v}^{2}}\ \overrightarrow{R}-\left( \overrightarrow{v}\cdot \overrightarrow{R} \right)\ \overrightarrow{v}}{cK}\]

Теперь ищем выражение для ротора запаздывающей скорости

\[rot\ \overrightarrow{v}=\begin{matrix}
   \overrightarrow{i}\left( \frac{\partial {{v}_{z}}}{\partial y}-\frac{\partial {{v}_{y}}}{\partial z} \right)  \\
   \overrightarrow{j}\left( \frac{\partial {{v}_{x}}}{\partial z}-\frac{\partial {{v}_{z}}}{\partial x} \right)  \\
   \overrightarrow{k}\left( \frac{\partial {{v}_{y}}}{\partial x}-\frac{\partial {{v}_{x}}}{\partial y} \right)  \\
\end{matrix}=\begin{matrix}
   \overrightarrow{i}\left( \frac{\partial {{v}_{z}}}{\partial t'}\frac{\partial t'}{\partial y}-\frac{\partial {{v}_{y}}}{\partial t'}\frac{\partial t'}{\partial z} \right)  \\
   \overrightarrow{j}\left( \frac{\partial {{v}_{x}}}{\partial t'}\frac{\partial t'}{\partial z}-\frac{\partial {{v}_{z}}}{\partial t'}\frac{\partial t'}{\partial x} \right)  \\
   \overrightarrow{k}\left( \frac{\partial {{v}_{y}}}{\partial t'}\frac{\partial t'}{\partial x}-\frac{\partial {{v}_{x}}}{\partial t'}\frac{\partial t'}{\partial y} \right)  \\
\end{matrix}=\begin{matrix}
   \frac{\overrightarrow{i}}{c}\left( -{{a}_{z}}\frac{\partial R}{\partial y}+{{a}_{y}}\frac{\partial R}{\partial z} \right)  \\
   \frac{\overrightarrow{j}}{c}\left( -{{a}_{x}}\frac{\partial R}{\partial z}+{{a}_{z}}\frac{\partial R}{\partial x} \right)  \\
   \frac{\overrightarrow{k}}{c}\left( -{{a}_{y}}\frac{\partial R}{\partial x}+{{a}_{x}}\frac{\partial R}{\partial y} \right)  \\
\end{matrix}=\begin{matrix}
   \frac{\overrightarrow{i}}{cK}\left( -{{a}_{z}}{{R}_{y}}+{{a}_{y}}{{R}_{z}} \right)  \\
   \frac{\overrightarrow{j}}{cK}\left( -{{a}_{x}}{{R}_{z}}+{{a}_{z}}{{R}_{x}} \right)  \\
   \frac{\overrightarrow{k}}{cK}\left( -{{a}_{y}}{{R}_{x}}+{{a}_{x}}{{R}_{y}} \right)  \\
\end{matrix}\]

 Далее векторное произведение радиус-вектора на ротор
запаздывающей скорости

\[cK\left( \overrightarrow{R}\times rot\ \overrightarrow{v} \right)=\left| \begin{matrix}
   \overrightarrow{i} & \overrightarrow{j} & \overrightarrow{k}  \\
   {{R}_{x}} & {{R}_{y}} & {{R}_{z}}  \\
   {{a}_{y}}{{R}_{z}}-{{a}_{z}}{{R}_{y}} & {{a}_{z}}{{R}_{x}}-{{a}_{x}}{{R}_{z}} & {{a}_{x}}{{R}_{y}}-{{a}_{y}}{{R}_{x}}  \\
\end{matrix} \right|=\begin{matrix}
   \overrightarrow{i}\left( {{R}_{y}}\left( {{a}_{x}}{{R}_{y}}-{{a}_{y}}{{R}_{x}} \right)-{{R}_{z}}\left( {{a}_{z}}{{R}_{x}}-{{a}_{x}}{{R}_{z}} \right) \right)  \\
   \overrightarrow{j}\left( {{R}_{z}}\left( {{a}_{y}}{{R}_{z}}-{{a}_{z}}{{R}_{y}} \right)-{{R}_{x}}\left( {{a}_{x}}{{R}_{y}}-{{a}_{y}}{{R}_{x}} \right) \right)  \\
   \overrightarrow{k}\left( {{R}_{x}}\left( {{a}_{z}}{{R}_{x}}-{{a}_{x}}{{R}_{z}} \right)-{{R}_{y}}\left( {{a}_{y}}{{R}_{z}}-{{a}_{z}}{{R}_{y}} \right) \right)  \\
\end{matrix}\] \[=\begin{matrix}
   \overrightarrow{i}\left( {{a}_{x}}\left( {{R}_{y}}^{2}+{{R}_{z}}^{2} \right)-{{R}_{x}}{{R}_{y}}{{a}_{y}}-{{R}_{x}}{{R}_{z}}{{a}_{z}} \right)  \\
   \overrightarrow{j}\left( {{a}_{y}}\left( {{R}_{x}}^{2}+{{R}_{z}}^{2} \right)-{{R}_{y}}{{R}_{x}}{{a}_{x}}-{{R}_{y}}{{R}_{z}}{{a}_{z}} \right)  \\
   \overrightarrow{k}\left( {{a}_{z}}\left( {{R}_{x}}^{2}+{{R}_{y}}^{2} \right)-{{R}_{z}}{{R}_{x}}{{a}_{x}}-{{R}_{z}}{{R}_{y}}{{a}_{y}} \right)  \\
\end{matrix}=\begin{matrix}
   \overrightarrow{i}\left( {{a}_{x}}\left( {{R}^{2}}-{{R}_{x}}^{2} \right)-{{R}_{x}}\left( {{a}_{y}}{{R}_{y}}+{{a}_{z}}{{R}_{z}} \right) \right)  \\
   \overrightarrow{j}\left( {{a}_{y}}\left( {{R}^{2}}-{{R}_{y}}^{2} \right)-{{R}_{y}}\left( {{a}_{x}}{{R}_{x}}+{{a}_{z}}{{R}_{z}} \right) \right)  \\
   \overrightarrow{k}\left( {{a}_{z}}\left( {{R}^{2}}-{{R}_{z}}^{2} \right)-{{R}_{z}}\left( {{a}_{x}}{{R}_{x}}+{{a}_{y}}{{R}_{y}} \right) \right)  \\
\end{matrix}=\begin{matrix}
   \overrightarrow{i}\left( {{a}_{x}}{{R}^{2}}-{{R}_{x}}\left( \overrightarrow{a}\cdot \overrightarrow{R} \right) \right)  \\
   \overrightarrow{j}\left( {{a}_{y}}{{R}^{2}}-{{R}_{y}}\left( \overrightarrow{a}\cdot \overrightarrow{R} \right) \right)  \\
   \overrightarrow{k}\left( {{a}_{z}}{{R}^{2}}-{{R}_{z}}\left( \overrightarrow{a}\cdot \overrightarrow{R} \right) \right)  \\
\end{matrix}\]
\[\overrightarrow{R}\times rot\ \overrightarrow{v}=\frac{{{R}^{2}}\overrightarrow{a}-\left( \overrightarrow{a}\cdot \overrightarrow{R} \right)\overrightarrow{R}}{cK}\]
Субстанционная производная радиус вектора
\[\left( \overrightarrow{v},\nabla  \right)\overrightarrow{R}={{v}_{x}}\frac{\partial }{\partial x}\overrightarrow{R}+{{v}_{y}}\frac{\partial }{\partial y}\overrightarrow{R}+{{v}_{z}}\frac{\partial }{\partial z}\overrightarrow{R}\]
Учитывая ранее найденные выражения производной радиус вектора по
координате

\[\frac{\partial \overrightarrow{R}}{\partial x}=\overrightarrow{i}+\frac{\overrightarrow{v}}{c}\frac{\partial R}{\partial x}\]

и если учесть что

\[\frac{\partial R}{\partial x}=\frac{{{R}_{x}}}{K}\]

то можно записать

\[\frac{\partial \overrightarrow{R}}{\partial x}=\overrightarrow{i}+\frac{\overrightarrow{v}}{c}\frac{{{R}_{x}}}{K}\]

\[\begin{align}
  & \left( \overrightarrow{v},\nabla  \right)\overrightarrow{R}={{v}_{x}}\left( \overrightarrow{i}+\frac{\overrightarrow{v}}{c}\frac{{{R}_{x}}}{K} \right)+{{v}_{y}}\left( \overrightarrow{j}+\frac{\overrightarrow{v}}{c}\frac{{{R}_{y}}}{K} \right)+{{v}_{z}}\left( \overrightarrow{k}+\frac{\overrightarrow{v}}{c}\frac{{{R}_{z}}}{K} \right)=\overrightarrow{v}+\frac{\overrightarrow{v}}{c}\frac{\left( \overrightarrow{v}\cdot \overrightarrow{R} \right)}{K}=\overrightarrow{v}\left( 1+\frac{\left( \overrightarrow{v}\cdot \overrightarrow{R} \right)}{cK} \right) \\ 
 & =\overrightarrow{v}\left( \frac{c\left( R-\frac{\overrightarrow{v}\cdot \overrightarrow{R}}{c} \right)}{cK}+\frac{\left( \overrightarrow{v}\cdot \overrightarrow{R} \right)}{cK} \right)=\overrightarrow{v}\left( \frac{cR-\overrightarrow{v}\cdot \overrightarrow{R}}{cK}+\frac{\left( \overrightarrow{v}\cdot \overrightarrow{R} \right)}{cK} \right)=\overrightarrow{v}\left( \frac{cR}{cK} \right)=\overrightarrow{v}\frac{R}{K}
\end{align}\]

Теперь

\[\left( \overrightarrow{R},\nabla  \right)\overrightarrow{v}={{R}_{x}}\frac{\partial }{\partial x}\overrightarrow{v}+{{R}_{y}}\frac{\partial }{\partial y}\overrightarrow{v}+{{R}_{z}}\frac{\partial }{\partial z}\overrightarrow{v}\]
\[\frac{\partial }{\partial x}\overrightarrow{v}=\frac{\partial \overrightarrow{v}}{\partial t'}\frac{\partial t'}{\partial x}=\overrightarrow{a}\frac{\partial t'}{\partial x}=-\frac{\overrightarrow{a}}{c}\frac{\partial R}{\partial x}=-\frac{\overrightarrow{a}}{c}\frac{{{R}_{x}}}{K}\]
\[\left( \overrightarrow{R},\nabla  \right)\overrightarrow{v}=-{{R}_{x}}\frac{\overrightarrow{a}}{c}\frac{{{R}_{x}}}{K}-{{R}_{y}}\frac{\overrightarrow{a}}{c}\frac{{{R}_{y}}}{K}-{{R}_{z}}\frac{\overrightarrow{a}}{c}\frac{{{R}_{z}}}{K}=-\frac{\overrightarrow{a}}{c}\frac{{{R}^{2}}}{K}\]

Таким образом, градиент скалярного произведения скорости на радиус
вектор

\[\begin{align}
  & \nabla \left( \overrightarrow{v}\cdot \overrightarrow{R} \right)=\overrightarrow{v}\times rot\overrightarrow{R}+\overrightarrow{R}\times rot\overrightarrow{v}+\left( \overrightarrow{v},\nabla  \right)\overrightarrow{R}+\left( \overrightarrow{R},\nabla  \right)\overrightarrow{v}= \\ 
 & =\left( \frac{{{v}^{2}}\ \overrightarrow{R}-\left( \overrightarrow{v}\cdot \overrightarrow{R} \right)\ \overrightarrow{v}}{cK} \right)+\left( \frac{{{R}^{2}}\overrightarrow{a}-\left( \overrightarrow{a}\cdot \overrightarrow{R} \right)\overrightarrow{R}}{cK} \right)+\left( \overrightarrow{v}+\frac{\overrightarrow{v}}{c}\frac{\left( \overrightarrow{v}\cdot \overrightarrow{R} \right)}{K} \right)+\left( -\frac{\overrightarrow{a}}{c}\frac{{{R}^{2}}}{K} \right) \\ 
 & =\left( \frac{{{v}^{2}}\ \overrightarrow{R}}{cK} \right)-\left( \frac{\left( \overrightarrow{a}\cdot \overrightarrow{R} \right)\overrightarrow{R}}{cK} \right)+\overrightarrow{v} \\ 
\end{align}\]

 Интересно, что этот градиент можно также найти значительно
более коротким путём - через сумму производных по координате

\[\nabla \left( \overrightarrow{v}\cdot \overrightarrow{R} \right)=\overrightarrow{i}\frac{\partial }{\partial x}\left( \overrightarrow{v}\cdot \overrightarrow{R} \right)+\overrightarrow{j}\frac{\partial }{\partial y}\left( \overrightarrow{v}\cdot \overrightarrow{R} \right)+\overrightarrow{k}\frac{\partial }{\partial z}\left( \overrightarrow{v}\cdot \overrightarrow{R} \right)\]
производная по координате
\[\frac{\partial }{\partial x}\left( \overrightarrow{v}\cdot \overrightarrow{R} \right)=\overrightarrow{v}\frac{\partial \overrightarrow{R}}{\partial x}+\overrightarrow{R}\frac{\partial \overrightarrow{v}}{\partial x}\]
\[\frac{\partial \overrightarrow{v}}{\partial x}=\frac{\partial \overrightarrow{v}}{\partial t'}\frac{\partial t'}{\partial x}=\overrightarrow{a}\left( \frac{\partial }{\partial x}\left\{ t-\frac{R}{c} \right\} \right)=\overrightarrow{a}\left( -\frac{1}{c}\frac{\partial R}{\partial x} \right)=-\frac{\overrightarrow{a}}{c}\frac{{{R}_{x}}}{K}\]
\[\frac{\partial }{\partial x}\left( \overrightarrow{v}\cdot \overrightarrow{R} \right)=\overrightarrow{v}\left( \overrightarrow{i}+\frac{\overrightarrow{v}}{c}\frac{{{R}_{x}}}{K} \right)+\overrightarrow{R}\left( -\frac{\overrightarrow{a}}{c}\frac{{{R}_{x}}}{K} \right)={{v}_{x}}+\left( \frac{{{v}^{2}}-\overrightarrow{R}\cdot \overrightarrow{a}}{c} \right)\frac{{{R}_{x}}}{K}\]
Складывая три компоненты приходим к тому же выражению
\[\nabla \left( \overrightarrow{v}\cdot \overrightarrow{R} \right)=\overrightarrow{v}+\left( \frac{{{v}^{2}}-\overrightarrow{R}\cdot \overrightarrow{a}}{c} \right)\frac{\overrightarrow{R}}{K}\]
Откуда
\[\nabla K=\nabla R-\frac{1}{c}\nabla \left( \overrightarrow{v}\cdot \overrightarrow{R} \right)=\left( \frac{\overrightarrow{R}\ }{K} \right)-\frac{1}{c}\left\{ \overrightarrow{v}+\left( \frac{{{v}^{2}}-\overrightarrow{R}\cdot \overrightarrow{a}}{c} \right)\frac{\overrightarrow{R}}{K} \right\}=\frac{\overrightarrow{R}\ }{K}\left( 1+\frac{\overrightarrow{a}\cdot \overrightarrow{R}-{{v}^{2}}}{{{c}^{2}}} \right)-\frac{\overrightarrow{v}}{c}\]
теперь
\[\nabla \frac{1}{\left( R-\frac{\overrightarrow{v}\cdot \overrightarrow{R}}{c} \right)}=-\frac{1}{{{\left( R-\frac{\overrightarrow{v}\cdot \overrightarrow{R}}{c} \right)}^{2}}}\left\{ \frac{\overrightarrow{R}}{\left( R-\frac{\overrightarrow{v}\cdot \overrightarrow{R}}{c} \right)}\left( 1+\frac{\overrightarrow{a}\cdot \overrightarrow{R}}{{{c}^{2}}}-\frac{\overrightarrow{v}\cdot \overrightarrow{v}}{{{c}^{2}}} \right)-\frac{\overrightarrow{v}}{c} \right\}\]
Откуда
\[\nabla \overrightarrow{{{A}_{H}}}=\frac{e}{c}\left\{ -\frac{\overrightarrow{a}\cdot \overrightarrow{R}}{c{{\left( R-\frac{\overrightarrow{v}\cdot \overrightarrow{R}}{c} \right)}^{2}}}-\frac{\overrightarrow{v}}{{{\left( R-\frac{\overrightarrow{v}\cdot \overrightarrow{R}}{c} \right)}^{2}}}\cdot \left\{ \frac{\overrightarrow{R}}{\left( R-\frac{\overrightarrow{v}\cdot \overrightarrow{R}}{c} \right)}\left( 1+\frac{\overrightarrow{a}\cdot \overrightarrow{R}}{{{c}^{2}}}-\frac{\overrightarrow{v}\cdot \overrightarrow{v}}{{{c}^{2}}} \right)-\frac{\overrightarrow{v}}{c} \right\} \right\}\]
На слагаемые
\[\nabla \overrightarrow{{{A}_{H}}}=\frac{e}{c}\left\{ -\frac{\overrightarrow{a}\cdot \overrightarrow{R}}{c{{\left( R-\frac{\overrightarrow{v}\cdot \overrightarrow{R}}{c} \right)}^{2}}}-\frac{\overrightarrow{v}\cdot \overrightarrow{R}}{{{\left( R-\frac{\overrightarrow{v}\cdot \overrightarrow{R}}{c} \right)}^{3}}}\left( 1+\frac{\overrightarrow{a}\cdot \overrightarrow{R}}{{{c}^{2}}}-\frac{\overrightarrow{v}\cdot \overrightarrow{v}}{{{c}^{2}}} \right)+\frac{\overrightarrow{v}\cdot \overrightarrow{v}}{{{\left( R-\frac{\overrightarrow{v}\cdot \overrightarrow{R}}{c} \right)}^{2}}} \right\}\]
\[\nabla \overrightarrow{{{A}_{H}}}=\frac{e}{c}\left\{ \frac{\overrightarrow{v}\cdot \overrightarrow{v}-\overrightarrow{a}\cdot \overrightarrow{R}}{c{{\left( R-\frac{\overrightarrow{v}\cdot \overrightarrow{R}}{c} \right)}^{2}}}-\frac{\overrightarrow{v}\cdot \overrightarrow{R}}{{{\left( R-\frac{\overrightarrow{v}\cdot \overrightarrow{R}}{c} \right)}^{3}}}\left( 1+\frac{\overrightarrow{a}\cdot \overrightarrow{R}}{{{c}^{2}}}-\frac{\overrightarrow{v}\cdot \overrightarrow{v}}{{{c}^{2}}} \right) \right\}\]
Теперь проверяем полученные выражения на соответствие калибровке Лоренца
\[div\ \overrightarrow{{{A}_{H}}}=-\frac{\varepsilon }{c}\frac{\partial \varphi }{\partial t}\]
Правая часть этого выражения
\[-\frac{\varepsilon }{c}\frac{\partial \varphi }{\partial t}=\frac{e}{}\frac{R\frac{\vec{v}\cdot \vec{v}-\overrightarrow{a}\cdot \overrightarrow{R}}{c}-~\overrightarrow{v}\cdot \overrightarrow{R}}{{{\left( R-\frac{\overrightarrow{v}\cdot \overrightarrow{R}}{c} \right)}^{3}}}\]
Итак, проверим тождество
\[\frac{\overrightarrow{v}\cdot \overrightarrow{v}-\overrightarrow{a}\cdot \overrightarrow{R}}{c{{\left( R-\frac{\overrightarrow{v}\cdot \overrightarrow{R}}{c} \right)}^{2}}}-\frac{\overrightarrow{v}\cdot \overrightarrow{R}}{{{\left( R-\frac{\overrightarrow{v}\cdot \overrightarrow{R}}{c} \right)}^{3}}}\left( 1+\frac{\overrightarrow{a}\cdot \overrightarrow{R}}{{{c}^{2}}}-\frac{\overrightarrow{v}\cdot \overrightarrow{v}}{{{c}^{2}}} \right)=\frac{R\frac{\vec{v}\cdot \vec{v}-\overrightarrow{a}\cdot \overrightarrow{R}}{c}-~\overrightarrow{v}\cdot \overrightarrow{R}}{{{\left( R-\frac{\overrightarrow{v}\cdot \overrightarrow{R}}{c} \right)}^{3}}}\]

\[\frac{\overrightarrow{v}\cdot \overrightarrow{v}-\overrightarrow{a}\cdot \overrightarrow{R}}{c{{\left( R-\frac{\overrightarrow{v}\cdot \overrightarrow{R}}{c} \right)}^{2}}}-\frac{\overrightarrow{v}\cdot \overrightarrow{R}}{{{\left( R-\frac{\overrightarrow{v}\cdot \overrightarrow{R}}{c} \right)}^{3}}}\left( 1+\frac{\overrightarrow{a}\cdot \overrightarrow{R}}{{{c}^{2}}}-\frac{\overrightarrow{v}\cdot \overrightarrow{v}}{{{c}^{2}}} \right)=\frac{R\left( \vec{v}\cdot \vec{v}-\overrightarrow{a}\cdot \overrightarrow{R} \right)}{{{\left( R-\frac{\overrightarrow{v}\cdot \overrightarrow{R}}{c} \right)}^{3}}}-\frac{~\overrightarrow{v}\cdot \overrightarrow{R}}{{{\left( R-\frac{\overrightarrow{v}\cdot \overrightarrow{R}}{c} \right)}^{3}}}\]

\[\frac{\overrightarrow{v}\cdot \overrightarrow{v}-\overrightarrow{a}\cdot \overrightarrow{R}}{c{{\left( R-\frac{\overrightarrow{v}\cdot \overrightarrow{R}}{c} \right)}^{2}}}-\frac{\overrightarrow{v}\cdot \overrightarrow{R}}{{{\left( R-\frac{\overrightarrow{v}\cdot \overrightarrow{R}}{c} \right)}^{3}}}\left( \frac{\overrightarrow{a}\cdot \overrightarrow{R}}{{{c}^{2}}}-\frac{\overrightarrow{v}\cdot \overrightarrow{v}}{{{c}^{2}}} \right)=\frac{R\left( \vec{v}\cdot \vec{v}-\overrightarrow{a}\cdot \overrightarrow{R} \right)}{{{\left( R-\frac{\overrightarrow{v}\cdot \overrightarrow{R}}{c} \right)}^{3}}}\]

\[\frac{\overrightarrow{v}\cdot \overrightarrow{v}-\overrightarrow{a}\cdot \overrightarrow{R}}{c{{\left( R-\frac{\overrightarrow{v}\cdot \overrightarrow{R}}{c} \right)}^{2}}}+\frac{\overrightarrow{v}\cdot \overrightarrow{R}}{{{\left( R-\frac{\overrightarrow{v}\cdot \overrightarrow{R}}{c} \right)}^{3}}}\left( \frac{\overrightarrow{v}\cdot \overrightarrow{v}-\overrightarrow{a}\cdot \overrightarrow{R}}{{{c}^{2}}} \right)=\frac{R\left( \vec{v}\cdot \vec{v}-\overrightarrow{a}\cdot \overrightarrow{R} \right)}{{{\left( R-\frac{\overrightarrow{v}\cdot \overrightarrow{R}}{c} \right)}^{3}}}\]
Сокращаем на
\[\frac{\overrightarrow{v}\cdot \overrightarrow{v}-\overrightarrow{a}\cdot \overrightarrow{R}}{{{c}^{2}}}\]
\[\frac{c}{{{\left( R-\frac{\overrightarrow{v}\cdot \overrightarrow{R}}{c} \right)}^{2}}}+\frac{\overrightarrow{v}\cdot \overrightarrow{R}}{{{\left( R-\frac{\overrightarrow{v}\cdot \overrightarrow{R}}{c} \right)}^{3}}}=\frac{cR}{{{\left( R-\frac{\overrightarrow{v}\cdot \overrightarrow{R}}{c} \right)}^{3}}}\]
Приводим к общему знаменателю
\[\frac{c\left( R-\frac{\overrightarrow{v}\cdot \overrightarrow{R}}{c} \right)}{{{\left( R-\frac{\overrightarrow{v}\cdot \overrightarrow{R}}{c} \right)}^{3}}}+\frac{\overrightarrow{v}\cdot \overrightarrow{R}}{{{\left( R-\frac{\overrightarrow{v}\cdot \overrightarrow{R}}{c} \right)}^{3}}}=\frac{cR}{{{\left( R-\frac{\overrightarrow{v}\cdot \overrightarrow{R}}{c} \right)}^{3}}}\]

\[\frac{cR-\overrightarrow{v}\cdot \overrightarrow{R}+\overrightarrow{v}\cdot \overrightarrow{R}}{{{\left( R-\frac{\overrightarrow{v}\cdot \overrightarrow{R}}{c} \right)}^{3}}}=\frac{cR}{{{\left( R-\frac{\overrightarrow{v}\cdot \overrightarrow{R}}{c} \right)}^{3}}}\]
Тождество доказано

    \begin{tcolorbox}[breakable, size=fbox, boxrule=1pt, pad at break*=1mm,colback=cellbackground, colframe=cellborder]
\prompt{In}{incolor}{ }{\boxspacing}
\begin{Verbatim}[commandchars=\\\{\}]

\end{Verbatim}
\end{tcolorbox}


    % Add a bibliography block to the postdoc
    
    
    
\end{document}
